\scanpage{316}

\mspage{55}
% archive sheet number: 315

de $*$-produits de $n$, resp. de $\sum_{i\in[1,n]}\ell_i$ facteurs de la forme $f_{\varepsilon,\unsure{\nu|n}}$. Pour terminer de tenir en mains l'objet \unsure{$1/2$ simpl.} tronqué $(\operatorname{Nerf}_i(\underline{\mathrm{Ty}})\unsure{]|)}_{0\leq i\leq2}$, il reste à décrire les \inserted{deux} applications de dégénérescence $\operatorname{Nerf}_1\unsure{\text{ --- }|\longrightarrow}\operatorname{Nerf}_2$ (cas où $f$ ou $g$ est une flèche identique), et surtout l'application de \uline{composition} des flèches
\[
\operatorname{Nerf}_2(\underline{\mathrm{Ty}})\xrightarrow{\gamma}\operatorname{Nerf}_1(\underline{\mathrm{Ty}})=\operatorname{Fl}(\mathrm{Ty}).
\]
Il suffit de donner $\gamma$ sur les générateurs, et on a \unsure{tautologiquement}, par définition des $f_{\varepsilon,\ell,\pi}$
\[
\gamma(f_{\varepsilon,\ell,\pi})=f_{\varepsilon,\operatorname{long}(\pi)}.
\]
(NB Il n'existe \inserted{au plus} qu'un seul morphisme de $\Delta_\varepsilon$ i.e. (de $\Delta_+$ ou de $\Delta_-$) dans un intervalle $I_\tau$. Il existe sss $\tau$ est de la forme $\Delta_\varepsilon^\ell$, $\ell\in\mathbb N$.)

Le cas $g=\mathrm{id}$ correspond au cas où
\[
\pi_i(j)=1\qquad\forall j\in[1,\ell_i],
\]
les $f_{\varepsilon_i,\ell_i,\pi_i}$ sont \unsure{donc|alors} déterminés en fonction des $f_{\varepsilon_i,\ell_i}$, et l'application de dégénérescence est
\[
 f_{\deleted{\unsure{+}}\varepsilon,\ell}\longmapsto f_{\varepsilon,\ell,\underline{1}_{\deleted{\unsure{n}}\ell}}
\]
\margin{(où $1_\ell$ est la fonction \unsure{conste|const.} de val. $1$ sur $[1,\ell]$.)}
\deleted{\unsure{du}} \unsure{Cas|Le} / $f=\mathrm{id}_{\unsure{x|y}}$ \inserted{est celui où} les $\ell_i$ sont égaux à $1$. Alors \deleted{les \unsure{fonctions}} \deleted{$\unsure{\pi_i}$} chaque $\pi_i:[1,1]\to\mathbb N$ peut être vu comme un \deleted{\unsure{enti}} élément de $\mathbb N$, la donnée des $(\varepsilon_i,\ell_i,\pi_i)$ revient à celle des $(\varepsilon_i,n_i)\in[\pm1]\times\mathbb N$, i.e. la donnée \deleted{\unsure{du système}}
